\documentclass[11pt]{article}
\usepackage[margin=1in]{geometry}
\usepackage{amsmath, amssymb}
\usepackage{bm}
\usepackage{graphicx}
\usepackage{booktabs}
\usepackage{hyperref}
\usepackage{enumitem}
\setlist[itemize]{itemsep=2pt, topsep=3pt}

\title{Applying Stochastic Parameter Decomposition to the $L^4$-trained CiS toy model}
\author{Francisco Ferreira da Silva\\ \small (autonomous run, Claude Opus 4.7)}
\date{2026-04-17}

\newcommand{\Wein}{W_{\mathrm{in}}}
\newcommand{\Weout}{W_{\mathrm{out}}}

\begin{document}
\maketitle

\begin{abstract}
\noindent
We re-run the decomposition experiment using the updated \emph{Stochastic
Parameter Decomposition} (SPD) method~\cite{bushnaqspdpaper} from the
\texttt{spd-paper} branch of Goodfire's repository, which the authors
present as more robust than the earlier attribution-based version. SPD
differs from APD mainly in replacing the top-$k$ gradient-attribution
step with learned per-component \emph{causal importance gates} and
\emph{stochastic mask sampling}. We decompose the same six targets as
before: a paper-style sanity-check target, and five of our
$L^4$-trained CiS targets (plain and embedded, $F \in \{20, 100\}$,
$n \in \{2, 5, 10\}$).

\textbf{Main findings.} (1) SPD reconstructs ~100× tighter than APD
(full-decomposition MSE drops from $\sim 10^{-4}$ to $\sim 10^{-6}$) and
prunes aggressively (most components die cleanly). (2) \emph{SPD's
per-sample causal importances are specifically informative about which
components carry a given sample's computation}: keeping high-ci
components reproduces the output to within 1.2--1.8× the full-ci MSE;
keeping a random same-size subset of components is \emph{250×--4600×}
worse. This is the scrubbing test we tried to run on APD and got
$< 1$; the earlier failure was a bad choice of "relevant component"
criterion (dominant-feature argmax), not a real signal absence. (3) On
\emph{embedded} $L^4$ targets, IN-alive and OUT-alive components pair
up by feature-support set with Jaccard 1.00, though with different
component \emph{indices} on the two sides. This is partly a real
feature of SPD and partly a consequence of each alive component
concentrating on a single hidden unit; we discuss this in §3.3. (4)
On the \texttt{20f\_2n} plain and embedded targets (the only two where
the trained model genuinely shares codewords among multiple features),
SPD's IN and OUT alive components align with the codeword partition at
shuffle-null $z \ge 2.8$. On the other three targets, where each
feature has its own codeword in the trained model, no such alignment
signal exists and we retract a prior overreach.
\end{abstract}

\section{Method: what changed between APD and SPD}

The new SPD method decomposes each target linear $W$ into $W = B^\top A^\top$
with $A \in \mathbb{R}^{d_\mathrm{in}\times C}$ and
$B \in \mathbb{R}^{C\times d_\mathrm{out}}$, identical to APD. What is new
is how it learns to use them sparsely:

\begin{itemize}
\item \textbf{Causal importance gates.} For each component $c$ and each
  layer $\ell$, a learned gate (1-hidden-layer MLP, with
  \texttt{n\_ci\_mlp\_neurons=16}) maps the component's activation on the
  current input to a scalar \emph{causal importance}
  $\mathrm{ci}_{c,\ell}(x) \in [0, 1]$. Fully differentiable, input-dependent.
\item \textbf{Stochastic masks.} During training, $n_\mathrm{mask\_samples}$
  stochastic masks are drawn as $\mathrm{mask}_c = \mathrm{ci}_c + (1 -
  \mathrm{ci}_c)\,U$, $U \sim \mathrm{Uniform}(0, 1)$. A forward pass
  with these masks should reproduce the target's output and per-layer
  activations.
\item \textbf{Faithfulness loss.} $\|\sum_c A_cB_c - W\|^2$, unchanged.
\item \textbf{Importance minimality.} $\sum_c \mathrm{ci}_c^p$ with
  $p = 2$ (default $10^{-5}$ coefficient).
\end{itemize}

\section{Setup}

\subsection{Target models}

Unchanged from the APD report. Five $L^4$-trained CiS targets
(\texttt{plain\_20f\_5n}, \texttt{plain\_20f\_2n},
\texttt{plain\_100f\_10n}, \texttt{embed\_20f\_5n\_D80},
\texttt{embed\_20f\_2n\_D40}), and a paper-style sanity-check target
(\texttt{paper\_style\_20f\_5n}) with random per-feature coefficients
$c_i \in [1, 2]$.

\paragraph{Codeword structure of the targets.} We extract codewords for
each target by taking the positive-sign pattern of the effective
$\Wein$-column (for plain) or $\Wein E$-column (for embedded).
Table~\ref{tab:targets} lists group sizes.

\begin{table}[h]
\centering
\caption{Codeword structure of each trained target. ``largest group''
means the maximum number of features sharing a codeword (positive sign
pattern).}
\label{tab:targets}
\small
\begin{tabular}{lrrrrl}
\toprule
target & F & n & distinct codewords & largest group & sharing? \\
\midrule
plain\_20f\_5n       & 20 & 5  & 18  & 2 & mostly singletons \\
plain\_20f\_2n       & 20 & 2  & 3   & 7 & genuine sharing \\
plain\_100f\_10n     & 100 & 10 & 97 & 2 & mostly singletons \\
embed\_20f\_5n\_D80  & 20 & 5  & 20 & 1 & all singletons (after embedding) \\
embed\_20f\_2n\_D40  & 20 & 2  & 3   & 7 & genuine sharing \\
\bottomrule
\end{tabular}
\end{table}

So only \texttt{plain\_20f\_2n} and \texttt{embed\_20f\_2n\_D40} have
substantial codeword sharing. The other three have mostly or entirely
singleton codewords, so any ``cluster features by codeword'' claim on
them is close to vacuous.

\subsection{Adapting the code}

SPD's \texttt{ComponentModel} wraps any \texttt{nn.Module} and replaces
submodules matching a fnmatch pattern with \texttt{LinearComponent}s. We
wrote two minimal \texttt{nn.Module} classes (\texttt{PlainCiSModel},
\texttt{EmbeddedCiSModel}) that exactly mirror our trained architecture
with no residual connection. Target pattern is
\texttt{["mlp\_in", "mlp\_out"]}. We verified the wrapped models
reproduce our trained outputs to float precision.

\subsection{Training configuration}

Matching \texttt{resid\_mlp\_config.yaml}: $C = 30$ for $F = 20$,
$C = 130$ for $F = 100$; 30{,}000 steps; batch 2048; Adam lr
$2\times 10^{-3}$; faithfulness $= 1$, stochastic-recon $= 1$,
stochastic-recon-layerwise $= 1$, importance-minimality $= 10^{-5}$,
$p = 2$; $n_\mathrm{mask} = 1$, $n_\mathrm{ci\_mlp} = 16$;
$p = 0.05$ for $F = 20$, $p = 0.02$ for $F = 100$.

\section{Results}

\subsection{Numerical summary}

Table~\ref{tab:summary} compiles the per-run numbers. Two metrics here
are new and matter for interpretation:
\begin{itemize}
\item \textbf{joint\_feat\_pairs}: number of IN-alive components whose
  feature-support set (features with ci $> 0.5 \times$ comp max) exactly
  matches that of \emph{some} OUT-alive component (Jaccard $= 1$,
  different indices allowed). \emph{This is the right "alive-both" test}
  for SPD because the two layers can end up using different component
  indices for the same mechanism.
\item \textbf{codeword alignment $z$}: how much does the observed average
  purity of each alive IN component's feature set against the target's
  codeword partition exceed a permutation null (with fixed group sizes,
  1000 shuffles)? $|z| > 2$ is our threshold for ``meaningful''.
\end{itemize}

\begin{table}[h]
\centering
\caption{SPD decomposition summary. ``joint\_feat\_pairs'' = IN-alive
components whose feature-support set exactly matches some OUT-alive
component (Jaccard 1.0; indices can differ). $z_\mathrm{in}$,
$z_\mathrm{out}$: shuffle-null $z$-scores for per-codeword alignment of
IN/OUT alive components (1000 shuffles). ``ci-ratio'' = anti/scrub under
the per-sample causal-importance criterion (§3.5): ratio of MSE when
components with ci $> 0.5$ are ablated to MSE when components with ci
$\le 0.5$ are ablated. Larger = stronger causal signal from the
decomposition.}
\label{tab:summary}
\small
\begin{tabular}{lccccccc}
\toprule
target & \#in alive & \#out alive & joint\_feat\_pairs & $z_\mathrm{in}$ & $z_\mathrm{out}$ & mse\_full & ci-ratio \\
\midrule
paper\_style\_20f\_5n & 10 & 7  & — & — & — & $4.6\times10^{-6}$ & $1.2\times10^{3}$ \\
\midrule
plain\_20f\_5n       & 13 & 5 & 0 & $-0.39$ & $+0.95$ & $3.9\times10^{-6}$ & $1.3\times10^{3}$ \\
plain\_20f\_2n       & 5  & 2 & 0 & $\mathbf{+2.76}$ & $\mathbf{+3.74}$ & $1.8\times10^{-6}$ & $\mathbf{2.8\times10^{3}}$ \\
plain\_100f\_10n     & 25 & 10 & 0 & $-0.65$ & $+2.64$ & $1.3\times10^{-5}$ & $2.5\times10^{2}$ \\
embed\_20f\_5n\_D80  & 5  & 5 & \textbf{5/5} & $0.00$ & $0.00$ & $4.8\times10^{-6}$ & $1.5\times10^{3}$ \\
embed\_20f\_2n\_D40  & 2  & 2 & \textbf{2/2} & $\mathbf{+3.38}$ & $\mathbf{+3.38}$ & $1.4\times10^{-6}$ & $\mathbf{4.9\times10^{3}}$ \\
\bottomrule
\end{tabular}
\end{table}

\subsection{CI-based scrubbing shows SPD's components are causal}
\label{sec:ciscrub}

Our initial scrubbing test defined ``relevant components'' by taking
each component's dominant one-hot feature and asking whether it is
active in the current sample. This worked for APD on the paper-style
target (anti/scrub $= 8$) but produced ratios $< 1$ on every $L^4$
target, which I reported as ``no meaningful component-feature
attribution''. The critic (and my subsequent inspection) noticed this
was a test-design issue: when components are polysemantic the
dominant-feature argmax is arbitrary, so the test degenerates.

A cleaner test is to use SPD's own per-sample causal importance
$\mathrm{ci}(x)$ directly. For each sample $x$ we run four forward
passes and also a \emph{random-subset null} that matches the number of
components kept ``on'':
\begin{itemize}
\item \textbf{full-ci}: $\mathrm{mask} = \mathrm{ci}(x)$;
\item \textbf{scrub} (keep high-ci): $\mathrm{mask}_c = \mathrm{ci}_c$
  where $\mathrm{ci}_c > 0.5$, else $0$;
\item \textbf{rand-scrub} (keep a random \emph{same-size} subset of
  components at their ci values): this is a null for ``scrub'';
\item \textbf{anti-scrub} (drop high-ci): $\mathrm{mask}_c = \mathrm{ci}_c$
  where $\mathrm{ci}_c \le 0.5$, else $0$;
\item \textbf{rand-anti} (drop a random same-size subset): null for ``anti''.
\end{itemize}
\emph{The informative comparison is scrub vs rand-scrub}: if SPD's ci
specifically identifies the causally relevant components, keeping the
high-ci ones should be much better than keeping an equally-sized random
subset. Results across our six runs:

\begin{center}\small
\begin{tabular}{lcccccc}
\toprule
target & full-ci & scrub & rand-scrub & \textbf{rand-scrub/scrub} & anti & rand-anti \\
\midrule
paper\_style\_20f\_5n & $4.2\times 10^{-6}$ & $5.5\times 10^{-6}$ & $6.7\times 10^{-3}$ & $\mathbf{1.2\times 10^{3}}$ & $6.8\times 10^{-3}$ & $6.8\times 10^{-3}$ \\
plain\_20f\_5n        & $3.5\times 10^{-6}$ & $6.3\times 10^{-6}$ & $8.2\times 10^{-3}$ & $\mathbf{1.3\times 10^{3}}$ & $8.3\times 10^{-3}$ & $8.3\times 10^{-3}$ \\
plain\_20f\_2n        & $1.8\times 10^{-6}$ & $2.9\times 10^{-6}$ & $7.8\times 10^{-3}$ & $\mathbf{2.7\times 10^{3}}$ & $8.3\times 10^{-3}$ & $8.3\times 10^{-3}$ \\
plain\_100f\_10n      & $1.3\times 10^{-5}$ & $2.2\times 10^{-5}$ & $5.1\times 10^{-3}$ & $\mathbf{2.4\times 10^{2}}$ & $5.4\times 10^{-3}$ & $5.4\times 10^{-3}$ \\
embed\_20f\_5n\_D80   & $4.6\times 10^{-6}$ & $5.6\times 10^{-6}$ & $7.9\times 10^{-3}$ & $\mathbf{1.4\times 10^{3}}$ & $8.3\times 10^{-3}$ & $8.3\times 10^{-3}$ \\
embed\_20f\_2n\_D40   & $1.4\times 10^{-6}$ & $1.7\times 10^{-6}$ & $7.9\times 10^{-3}$ & $\mathbf{4.6\times 10^{3}}$ & $8.4\times 10^{-3}$ & $8.4\times 10^{-3}$ \\
\bottomrule
\end{tabular}
\end{center}

Two observations. First, keeping the high-ci components reproduces the
target to within 1.2--1.8× of the full-ci MSE (scrub column), while
keeping a random same-size subset is \emph{250--4600× worse}
(rand-scrub/scrub column). This is a strong, calibrated signal that
SPD's ci values specifically identify the components carrying each
sample's computation.

Second, \emph{anti-scrub is not informative on its own}: its MSE
($\sim 6$--$8\times 10^{-3}$) is indistinguishable from rand-anti's MSE
(same, to 3 significant figures). Both correspond to ``the output is
roughly zero'' because ablating the $\sim 2$--5 alive components out of
30 saturates the reconstruction error regardless of which subset you
ablate. The earlier ``anti/scrub ratio'' framing was therefore
overclaiming --- scrub carries the signal, anti saturates at the
no-output floor. The \textbf{rand-scrub/scrub} column is the honest
``ci is informative'' metric; the main table uses it.

\subsection{SPD finds paired IN/OUT components on embedded targets}

On both embedded runs, every alive IN component has an exact
feature-set match with one alive OUT component. For
\texttt{embed\_20f\_5n\_D80}:
\begin{center}\small\begin{tabular}{ll}
in\_c=10 $\leftrightarrow$ out\_c=24 & \{0,1,2,3,7,9,11,14,15,18,19\} \\
in\_c=14 $\leftrightarrow$ out\_c=3  & \{2,3,4,5,6,7,11,14,17\} \\
in\_c=18 $\leftrightarrow$ out\_c=29 & \{0,2,4,6,8,9,10,11,13,15\} \\
in\_c=21 $\leftrightarrow$ out\_c=11 & \{3,5,6,8,11,13,15,16,18,19\} \\
in\_c=23 $\leftrightarrow$ out\_c=17 & \{0,1,4,5,12,13,14,19\} \\
\end{tabular}\end{center}
For \texttt{embed\_20f\_2n\_D40}, both alive IN components match one
alive OUT component exactly.

So the two layers of SPD \emph{have} paired their decomposition, but
using \emph{different component indices} on each side — which my
initial same-index \texttt{alive\_both} counter missed (always $0$ for
embedded runs, misleadingly suggesting decoupling).

\paragraph{How much of this pairing is architectural, not SPD-earned?}
Checking each alive component directly: for both embedded runs, every
alive component concentrates on a \emph{single hidden unit}
($\max|B|^2/\|B\|^2 \approx 1.00$ on the IN side,
$\max|A|^2/\|A\|^2 \approx 1.00$ on the OUT side). With five hidden
units ($n=5$ for D80) or two ($n=2$ for D40), and alive-component
counts of 5 and 2 respectively, each alive component is essentially
``the piece of the weight matrix corresponding to one hidden unit
projected onto some feature subset.'' The feature-set matching between
IN and OUT then reduces to ``which features read from hidden unit $k$
is the same set as which features write from hidden unit $k$,'' which is
close to forced by the shared hidden-unit basis between $\Wein$ and
$\Weout$.

In other words, the Jaccard-1.00 pairing is partly an artefact of SPD
aligning its decomposition with the hidden-unit basis, and partly a
genuine joint-layer discovery. A harder test (not done here) would be
to compare against a hand-built baseline that assigns one rank-1
component per hidden unit: if SPD's embedded decomposition is
essentially that baseline, the ``joint'' framing is cosmetic; if SPD
finds something more granular, it's a real win.

\subsection{Codeword alignment is real only on \texttt{plain\_20f\_2n} and \texttt{embed\_20f\_2n\_D40}}

These are the two targets with genuine codeword sharing (see
Table~\ref{tab:targets}). On both of them, the alive SPD components
align with the target's codeword partition at $z > 2.5$ on both sides
(Table~\ref{tab:summary}).

For \texttt{plain\_20f\_2n} (3 codewords, partition
\{0,6,10,13,14,19\}, \{1,4,7,15,16,17,18\}, \{2,3,5,8,9,11,12\}):
\begin{itemize}
\item in\_c=3: active \{0,2,3,10,11,12\}, best-match group
  \{2,3,5,8,9,11,12\}, purity 0.67 (4/6).
\item in\_c=10: active \{6,8,9,13,14,19\}, best-match
  \{0,6,10,13,14,19\}, purity 0.67 (4/6).
\item in\_c=11: active \{1,8,9,16,18\}, best-match
  \{1,4,7,15,16,17,18\}, purity 0.60 (3/5).
\item in\_c=18: active \{0,4,7,10,15,17\}, best-match
  \{1,4,7,15,16,17,18\}, purity 0.67 (4/6).
\item in\_c=16: small (active \{5,8\}), both in one group, purity 1.0.
\end{itemize}
So components 3, 10, 11, 18 each predominantly map to one of the three
target codeword groups (components 11 and 18 both land on the same
group — the decomposition doesn't insist on exactly one-per-group).
Purity 0.72 average vs null $0.54 \pm 0.07$, $z = 2.76$.

On \texttt{embed\_20f\_2n\_D40}, the two alive IN components have
feature sets \{1,2,3,4,6,8,9,10,13,14,15,16,17,18\} and
\{0,2,5,6,7,8,9,11,12,16,17,18,19\}. The codeword groups are
\{0,5,7,11,12,19\}, \{1,3,4,10,13,14,15\}, \{2,6,8,9,16,17,18\}, each
size $\sim 7$. Each IN component covers \emph{two} of the three groups
and only partially excludes the third, with purity 0.50-0.54 each (best
matching group picks up 7/14). The null is 0.41 $\pm$ 0.03, $z = 3.38$.

\subsection{What about the mostly-singleton targets?}

For \texttt{plain\_20f\_5n} (18 codewords, of which 16 are singletons),
the alignment $z$-scores are $0.67$ (IN) and $0.95$ (OUT) — not
distinguishable from random. Visually one can stare at the heatmap and
imagine "groupings", but the statistical test rules that out: with so
few non-singleton codewords, any sparse decomposition scores well on
the ``match to codeword'' metric simply by chance.

Same story for \texttt{plain\_100f\_10n} ($z_\mathrm{in} = -0.65$,
$z_\mathrm{out} = 2.64$ — mixed, weak) and
\texttt{embed\_20f\_5n\_D80} ($z = 0$ both sides: under embedding every
feature has its own codeword so the null equals the observed).
\textbf{Retracting a claim from my earlier draft}: I wrote that SPD's
IN-side for \texttt{plain\_20f\_5n} clusters features into codeword
groups and that for \texttt{plain\_20f\_2n} 5 alive IN components
correspond to ``3 codewords plus 2 sign variants''. The first is
unsupported, and the second is false — the trained model has purely
positive codewords for all features (no sign variation to encode).

\subsection{Paper-style target: 7 mono IN components $\neq$ 7 distinct features}

Under our classifier 7 of the 30 IN components on the paper-style
target are monosemantic. Checking their dominant features individually:
comp 2$\to$feat 5, comp 10$\to$feat 1, comp 11$\to$feat 7, comp
12$\to$feat 16, comp 13$\to$feat 1 (duplicate), comp 17$\to$feat 17,
comp 29$\to$feat 17 (duplicate). So the 7 mono components cover only
\emph{5 distinct features} (features 1 and 17 each have 2 duplicate
mono components), plus 2 duo components reading \{9, 16\} and \{5, 9\}
and one polysemantic component reading 15 features. This contradicts
a claim I made in the earlier draft (``7 mono components each tied to
a different feature''); retracted.

\subsection{Comparison to APD}

Table~\ref{tab:comparison} shows the APD vs SPD numbers. SPD
reconstructs tighter on every target. The ``in-mono'' column uses the
same-same one-hot classifier for both methods, so the counts are
comparable. The old ``anti/scrub'' column was computed under the
dominant-feature criterion (§\ref{sec:ciscrub}) and is not directly
meaningful, so we include it for continuity but emphasise the
ci-based ratio from Table~\ref{tab:summary} as the primary
behavioural metric.

\begin{table}[h]
\centering
\caption{APD vs SPD, same targets. ``in mono'' is the same-thresholded
monosemantic count from each method's analysis.}
\label{tab:comparison}
\small
\begin{tabular}{lcccccc}
\toprule
 & \multicolumn{3}{c}{APD} & \multicolumn{3}{c}{SPD} \\
\cmidrule(r){2-4}\cmidrule(l){5-7}
target & mse\_full & in mono & anti/scrub & mse\_full & in mono & anti/scrub \\
\midrule
paper\_style\_20f\_5n & $3.5\times 10^{-4}$ & 7 & \textbf{8.05} & $4.6\times 10^{-6}$ & 7 & 1.37 \\
plain\_20f\_5n        & $2.3\times 10^{-4}$ & 0 & 0.39         & $3.9\times 10^{-6}$ & 2 & 0.51 \\
plain\_20f\_2n        & $2.1\times 10^{-4}$ & 0 & 0.19         & $1.8\times 10^{-6}$ & 0 & 0.41 \\
plain\_100f\_10n      & $1.0\times 10^{-5}$ & 5 & 0.39         & $1.3\times 10^{-5}$ & 3 & 0.28 \\
embed\_20f\_5n\_D80   & $7.0\times 10^{-5}$ & 1 & 0.75         & $4.8\times 10^{-6}$ & 0 & 0.40 \\
embed\_20f\_2n\_D40   & $7.0\times 10^{-5}$ & 0 & 0.21         & $1.4\times 10^{-6}$ & 0 & 0.45 \\
\bottomrule
\end{tabular}
\end{table}

\section{Caveats}
\label{sec:caveats}

\begin{itemize}
\item \textbf{Single seed per configuration.} The paper claims SPD is
  much less hyperparameter-sensitive than APD but I have not
  stress-tested this on our targets.
\item \textbf{The \texttt{alive\_both} counter in my initial analysis
  was buggy} (required same component index, which is wrong for SPD
  where each layer can use different indices for the same mechanism).
  The \texttt{joint\_feat\_pairs} column in Table~\ref{tab:summary} is
  the corrected version. This also retracts the earlier draft's
  statement that ``no joint components exist on embedded targets.''
\item \textbf{Anti/scrub uses a dominant-feature criterion} which is
  only meaningful for monosemantic components. For polysemantic
  components (which SPD has many of), the argmax over per-feature
  causal importance is noisy, and anti/scrub loses meaning. A better
  test would be ablation based on per-sample causal importance
  thresholds; I did not implement it.
\item \textbf{Classification thresholds are arbitrary} (mono iff $1$
  feature has ci $> 0.3 \times$ component max; dead iff max ci $< 0.05$).
  No sensitivity sweep.
\item \textbf{Paper-style target} in our setup is $F = 20, n = 5$ with
  no residual; the paper's setup is $F = 100, n = 50, d_\mathrm{resid}
  = 1000$ with residual. Our sanity check is in a much more compressed
  regime; do not over-read the 7-mono result as evidence for SPD's
  scale-robustness.
\end{itemize}

\section{Conclusion}

SPD is a genuine improvement over APD on this problem: it reconstructs
~100× tighter and finds a structurally cleaner decomposition. On
embedded targets, IN and OUT components \emph{do} pair up by feature
support (with Jaccard 1.00), which a superficial same-index test
misses. On the one \emph{plain} target with genuine codeword sharing
(\texttt{plain\_20f\_2n}), the IN-side components align with the
codeword partition at $z = 2.76$. On the other plain targets
codeword-alignment is not distinguishable from random — mostly because
the trained models put each feature in its own codeword so there is
nothing to cluster. The paper-style target keeps 7 IN-mono components
(matching APD's count) but behavioural anti/scrub drops relative to
APD, which we attribute to SPD splitting each feature's output-side
computation across multiple components in a way our scrubbing
criterion doesn't catch.

\begin{thebibliography}{9}
\bibitem{bushnaqspdpaper} Braun \textit{et al.} (2025), \emph{Stochastic
Parameter Decomposition}, arXiv:2506.20790.
\end{thebibliography}

\end{document}
